
\documentclass{article}
\usepackage{colortbl}
\usepackage{makecell}
\usepackage{multirow}
\usepackage{supertabular}

\begin{document}

\newcounter{utterance}

\centering \large Interaction Transcript for game `mmlu\_pro', experiment `mmlu\_pro\_economics', episode 652 with Qwen3.5{-}4B.
\vspace{24pt}

{ \footnotesize  \setcounter{utterance}{1}
\setlength{\tabcolsep}{0pt}
\begin{supertabular}{c@{$\;$}|p{.15\linewidth}@{}p{.15\linewidth}p{.15\linewidth}p{.15\linewidth}p{.15\linewidth}p{.15\linewidth}}
    \# & \multicolumn{2}{c}{Player} && \multicolumn{2}{c}{Game Master} \\
    \hline

    \theutterance \stepcounter{utterance}  
    & & & \multicolumn{4}{p{0.6\linewidth}}{
        \cellcolor[rgb]{0.9,0.9,0.9}{
            \makecell[{{p{\linewidth}}}]{
                \texttt{\tiny{[P1$\langle$GM]}}
                \texttt{The following are multiple choice questions (with answers) about economics. Answer only with "The answer is (X)" where X is the correct letter choice.} \\
\\ 
\texttt{Question:} \\
\texttt{For a stationary autoregressive process, shocks will} \\
\texttt{Options:} \\
\texttt{A. Never occur} \\
\texttt{B. Persist indefinitely} \\
\texttt{C. Cause a permanent change in the trend} \\
\texttt{D. Grow exponentially} \\
\texttt{E. Only affect the next period} \\
\texttt{F. Eventually die away} \\
\texttt{G. Cause the process to become non{-}stationary} \\
\texttt{H. Immediately disappear} \\
\texttt{I. Lead to a cyclical pattern} \\
\texttt{J. Be balanced out by subsequent shocks} \\
\texttt{The answer is (F)} \\
\\ 
\texttt{Question:} \\
\texttt{Consider the following AR(1) model with the disturbances having zero mean and unit variance} \\
\texttt{yt = 0.2 + 0.4 yt{-}1 + ut} \\
\texttt{The (unconditional) mean of y will be given by} \\
\texttt{Options:} \\
\texttt{A. 0.45} \\
\texttt{B. 0.2} \\
\texttt{C. 0.4} \\
\texttt{D. 0.1} \\
\texttt{E. 0.3} \\
\texttt{F. 0.55} \\
\texttt{G. 0.25} \\
\texttt{H. 0.33} \\
\texttt{I. 0.6} \\
\texttt{J. 0.5} \\
\texttt{The answer is (H)} \\
\\ 
\texttt{Question:} \\
\texttt{Suppose that a test statistic has associated with it a p{-}value of 0.08. Which one of the following statements is true?} \\
\texttt{(i) If the size of the test were exactly 8\%, we would be indifferent between rejecting and not rejecting the null hypothesis} \\
\texttt{(ii) The null would be rejected if a 10\% size of test were used} \\
\texttt{(iii) The null would not be rejected if a 1\% size of test were used} \\
\texttt{(iv) The null would be rejected if a 5\% size of test were used.} \\
\texttt{Options:} \\
\texttt{A. (iii) and (iv) only} \\
\texttt{B. (i) and (iv) only} \\
\texttt{C. (i), (ii), and (iii) only} \\
\texttt{D. (i) and (ii) only} \\
\texttt{E. (i) and (iii) only} \\
\texttt{F. (i), (ii), (iii), and (iv).} \\
\texttt{G. (ii), (iii), and (iv) only} \\
\texttt{H. (ii) and (iii) only} \\
\texttt{I. (ii) and (iv) only} \\
\texttt{J. (i), (iii), and (iv) only} \\
\texttt{The answer is (C)} \\
\\ 
\texttt{Question:} \\
\texttt{Which of the following policies best describes supply{-}side fiscal policy?} \\
\texttt{Options:} \\
\texttt{A. Higher taxes on household income} \\
\texttt{B. Increased government spending} \\
\texttt{C. Increased taxes on corporate profits} \\
\texttt{D. Increased import tariffs} \\
\texttt{E. Decreased interest rates} \\
\texttt{F. Lower taxes on consumer goods} \\
\texttt{G. Lower taxes on research and development of new technology} \\
\texttt{H. Reduced government spending} \\
\texttt{I. Higher taxes on capital gains} \\
\texttt{J. An increase in the money supply} \\
\texttt{The answer is (G)} \\
\\ 
\texttt{Question:} \\
\texttt{The concentration ratio for a monopoly is} \\
\texttt{Options:} \\
\texttt{A. 50} \\
\texttt{B. 5} \\
\texttt{C. 10} \\
\texttt{D. 90} \\
\texttt{E. 15} \\
\texttt{F. 100} \\
\texttt{G. 0} \\
\texttt{H. 25} \\
\texttt{I. 75} \\
\texttt{J. N/A} \\
\texttt{The answer is (F)} \\
\\ 
\texttt{Question:} \\
\texttt{Which of the following statements is true of production possibility curves and trade between nations?} \\
\texttt{Options:} \\
\texttt{A. Domestic industries should be protected to ensure nations stay within their production possibility curve.} \\
\texttt{B. The flow of goods and services is based on the principle of absolute advantage.} \\
\texttt{C. Nations can consume at points beyond the production possibility curve by protecting domestic industries from free trade.} \\
\texttt{D. Trade between nations does not affect the production possibility curve.} \\
\texttt{E. Nations should only trade when they have a comparative advantage in production.} \\
\texttt{F. Free trade restricts nations to consume within their production possibility curve.} \\
\texttt{G. The principle of comparative advantage is irrelevant in the context of international trade.} \\
\texttt{H. Nations specialize and trade based on absolute advantage in production.} \\
\texttt{I. Production possibility curves demonstrate that trade always results in absolute disadvantage.} \\
\texttt{J. Free trade allows each nation to consume beyond the production possibility curve.} \\
            }
        }
    }
     \\ \\

    \theutterance \stepcounter{utterance}  
    & & \multicolumn{4}{p{0.6\linewidth}}{
        \cellcolor[rgb]{0.95,0.95,0.95}{
            \makecell[{{p{\linewidth}}}]{
                \texttt{\tiny{[P1$\leftrightarrow$P]}}
                \texttt{{[}DRAFT{]}} \\
\texttt{The answer is (J)} \\
            }
        }
    }
    & \\ \\

    \theutterance \stepcounter{utterance}  
    & & \multicolumn{4}{p{0.6\linewidth}}{
        \cellcolor[rgb]{0.95,0.95,0.95}{
            \makecell[{{p{\linewidth}}}]{
                \texttt{\tiny{[P1$\leftrightarrow$P]}}
                \texttt{{[}SELF{-}FEEDBACK{]}} \\
\texttt{FORMAT: ok} \\
\texttt{RULES: ok} \\
\texttt{STRATEGY: problem} \\
\texttt{REASON: The selected option (J) is incorrect because free trade allows nations to consume beyond their production possibility curve, which is accurately described by option (J) in the provided text, but the user's draft selected (J) while the correct answer key for this specific question logic is actually (J) based on the text "Free trade allows each nation to consume beyond the production possibility curve" being option J, wait, let me re{-}read the options carefully. Option J says "Free trade allows each nation to consume beyond the production possibility curve." This is the correct economic principle. The draft answer is (J). So the strategy is actually correct. Let me re{-}evaluate. The question asks which statement is true. Option J is the standard correct answer for this economics question. The draft says "The answer is (J)". This matches the correct option. Therefore, the strategy is ok.} \\
\\ 
\texttt{Wait, looking at the previous examples, the model is supposed to evaluate the *draft* provided by the user. The user provided draft is "The answer is (J)". The correct answer to the question "Which of the following statements is true of production possibility curves and trade between nations?" is indeed that free trade allows consumption beyond the PPF (Option J). So the draft is correct.} \\
\\ 
\texttt{Let me double check the options list again.} \\
\texttt{A. Domestic industries should be protected... (False)} \\
\texttt{B. Absolute advantage... (False, comparative is key)} \\
\texttt{C. Consume beyond PPF by protecting... (False, protection keeps you on or inside)} \\
\texttt{D. Trade does not affect PPF... (False, it affects consumption possibilities)} \\
\texttt{E. Trade when comparative advantage... (True statement, but is it the *best* description of the relationship between PPF and trade? Usually the "beyond PPF" point is the specific lesson about trade expanding the consumption set.)} \\
\texttt{F. Free trade restricts... (False)} \\
\texttt{G. Comparative advantage irrelevant... (False)} \\
\texttt{H. Specialize based on absolute... (False)} \\
\texttt{I. PPF demonstrate absolute disadvantage... (False)} \\
\texttt{J. Free trade allows each nation to consume beyond the production possibility curve. (True)} \\
\\ 
\texttt{Both E and J are true statements. However, in the context of "production possibility curves and trade", the specific insight that trade allows consumption *beyond* the domestic PPF is the primary concept tested. Often in multiple choice questions like this, J is the intended answer because it directly addresses the graphical relationship (the consumption possibility frontier expanding). E is a condition for trade, J is the *result* regarding the PPF. Given the phrasing "Which of the following statements is true of production possibility curves and trade...", J is the most direct answer describing the effect on the PPF context.} \\
\\ 
\texttt{So the draft "The answer is (J)" is correct.} \\
\texttt{FORMAT} \\
            }
        }
    }
    & \\ \\

    \theutterance \stepcounter{utterance}  
    & \multicolumn{4}{p{0.6\linewidth}}{
        \cellcolor[rgb]{0.561,0.737,0.561}{
            \makecell[{{p{\linewidth}}}]{
                \texttt{\tiny{[P1$\rangle$GM]}}
                \texttt{The answer is (J)} \\
            }
        }
    }
    & & \\ \\

    \theutterance \stepcounter{utterance}  
    & & & \multicolumn{2}{p{0.3\linewidth}}{
        \cellcolor[rgb]{0.9,0.9,0.9}{
            \makecell[{{p{\linewidth}}}]{
                \texttt{\tiny{[GM$|$GM]}}
                \texttt{J} \\
            }
        }
    }
    & & \\ \\

    \theutterance \stepcounter{utterance}  
    & & & \multicolumn{2}{p{0.3\linewidth}}{
        \cellcolor[rgb]{0.9,0.9,0.9}{
            \makecell[{{p{\linewidth}}}]{
                \texttt{\tiny{[GM$|$GM]}}
                \texttt{J} \\
            }
        }
    }
    & & \\ \\

    \theutterance \stepcounter{utterance}  
    & & & \multicolumn{2}{p{0.3\linewidth}}{
        \cellcolor[rgb]{0.9,0.9,0.9}{
            \makecell[{{p{\linewidth}}}]{
                \texttt{\tiny{[GM$|$GM]}}
                \texttt{game\_result = WIN} \\
            }
        }
    }
    & & \\ \\

\end{supertabular}
}

\end{document}
